\chapter{Background}
\label{ch:background}

This chapter builds up the objects the rest of the report uses, starting from
plain definitions. Section~\ref{sec:bg-loop} defines the repair loop precisely.
Section~\ref{sec:bg-cost} explains what validation costs, why the order in
which tests run changes that cost, and why counting test runs is not the same as
measuring time. Section~\ref{sec:bg-hypothesis} explains what a repair
hypothesis is and why putting the whole history in the prompt does not stop the
model repeating one.

\section{Iterative LLM-Based Repair}
\label{sec:bg-loop}

\subsection{The Objects}

Table~\ref{tab:terms} lists every term this report uses in a technical sense.
Each is defined once here and used consistently afterwards.

\begin{table}[htbp]
\centering
\caption{Terms used in this report.}
\label{tab:terms}
\footnotesize
\begin{tabularx}{\textwidth}{@{}>{\raggedright\arraybackslash}p{0.20\textwidth} X@{}}
\toprule
Term & Meaning in this report\\
\midrule
Faulty program $s_0$ & The program that contains the bug.\\
Test case $x$ & One input together with the output the program is expected to produce for it.\\
Test pool $P$ & The finite set of test cases available for one task.\\
Reject (refute) & A test case \emph{rejects} a program when its output on that input differs from the expected one. One rejecting test proves the program wrong.\\
Candidate patch $s$ & A modified program produced by the LLM in one round.\\
Round & One pass of generate $\rightarrow$ validate $\rightarrow$ refine. A \emph{session} is all rounds spent on one task, up to the budget $B$.\\
Scan & One validation: running a candidate against test cases from $P$ until one rejects it, or until all have passed.\\
Plausible patch & A candidate no test case in $P$ rejects --- ``plausible'' rather than ``correct'' because passing a finite pool is weaker than being right~\citep{qi2015plausibility}.\\
Test-run cost $c(x)$ & The CPU time consumed by running one candidate on one test case $x$.\\
Repair hypothesis & The idea behind a patch: what the model believes is wrong and how it intends to fix it. Many patches can express one hypothesis.\\
\bottomrule
\end{tabularx}
\end{table}

\subsection{The Loop as a Formal Object}

\begin{definition}[Repair session]
\label{def:session}
A repair session with budget $B$ on a faulty program $s_0$ and a test pool $P$
is the following sequence, for $t = 1,2,\dots,B$:
\begin{enumerate}[leftmargin=*,itemsep=2pt,topsep=3pt]
  \item The LLM draws a candidate patch
        $s_t \sim \pi(\cdot \mid \mathcal{H}_{t-1})$, where $\mathcal{H}_{t-1}$
        is whatever the loop has chosen to put in the prompt from rounds
        $1,\dots,t-1$.
  \item The \emph{validator} runs $s_t$ against test cases from $P$ in some
        order $\sigma_t$, stopping at the first test case that rejects $s_t$.
  \item If no test case rejects $s_t$, the session \emph{accepts} $s_t$ and
        stops. Otherwise the rejecting test case, together with whatever else
        the loop chooses to carry, becomes part of $\mathcal{H}_t$, and the next
        round begins.
\end{enumerate}
\end{definition}

Two choices are left open by Definition~\ref{def:session}, and they are exactly
the two levers of this work. The order $\sigma_t$ is free, and nothing in the
definition says it must be the order the benchmark happens to ship. The contents
of $\mathcal{H}_t$ are free, and nothing says they must be the full transcript.

\subsection{The Two Bills}

A repair session spends two different resources, and keeping them apart matters
because optimising one can increase the other.

\begin{itemize}[leftmargin=*,itemsep=3pt]
\item \textbf{The generation bill}, paid in model tokens and therefore in money
or in GPU time. It grows with the number of rounds and with the size of the
prompt.
\item \textbf{The validation bill}, paid in CPU seconds spent executing
candidate programs. It grows with the number of rounds, with how many test cases
each scan runs, and with how expensive those particular test cases are.
\end{itemize}

This report targets the validation bill with its scheduler and the generation
bill with its hints, and reports both.

\section{Validation Cost and Test Ordering}
\label{sec:bg-cost}

\subsection{What a Session Costs}

Write $c(x)$ for the CPU cost of running one candidate on test case $x$, and
$\sigma_t$ for the order the validator uses in round $t$, so that
$\sigma_t(j)$ is the $j$-th test case it runs. The validation cost of a session
is
\begin{equation}
C \;=\; \sum_{t=1}^{T} \; \sum_{j=1}^{r_t} c\!\left(\sigma_t(j)\right),
\qquad
r_t \;=\; \min\{\, j \;:\; \sigma_t(j) \text{ rejects } s_t \,\},
\label{eq:cost}
\end{equation}
where $T$ is the number of rounds actually run and $r_t = |P|$ when nothing
rejects $s_t$. Two observations about Equation~\eqref{eq:cost} drive the whole
design.

First, $\sigma_t$ is the only term the validator controls. The candidate comes
from the model, the pool comes from the benchmark, and the per-test costs are
properties of the programs; the order is a free choice made on every round.

Second, minimising the \emph{number} of test runs $\sum_t r_t$ is the same
objective as minimising $C$ \textbf{only if $c$ is constant}. When test costs
differ, a schedule can run more tests and still finish sooner, or run fewer
tests and take longer. The two quantities can therefore rank two schedules in
opposite orders.

\subsection{Why Executions and Seconds Are Not Interchangeable}

The test-prioritisation literature usually reports the number of test
executions rather than elapsed time, and does so for a defensible reason. The
largest study of patch validation to date measures executions precisely because
elapsed time is unstable across machines and configurations, where an execution
count reproduces~\citep{lou2024aprregression}. That choice is sound when
executions are interchangeable --- when running any one test costs about as
much as running any other.

It stops being sound when they are not. In the setting studied here, a test case
is one whole-program execution on one input, and inputs within a single task's
pool differ by orders of magnitude in size and in the amount of work they
provoke. A deployment pays for seconds, not for executions, so a result reported
only in executions does not tell the deployment what it will pay. This is why
Chapter~\ref{ch:experiments} makes CPU seconds the primary unit and reports
executions alongside rather than instead.

\subsection{Order Changes the Cost but Not the Answer}

\begin{definition}[Soundness of a schedule]
\label{def:soundness}
A candidate is accepted only after surviving \emph{every} usable test case in
$P$. Acceptance is therefore independent of the order $\sigma_t$: no ordering
can accept a patch that another ordering rejects, and no ordering can reject a
patch that another ordering accepts.
\end{definition}

Definition~\ref{def:soundness} is what makes a pure cost claim possible. If two
schedules repaired different numbers of bugs, the difference would be an
implementation bug rather than a finding, so the experiment asserts the
equality rather than reporting it (Section~\ref{sec:exp-metrics}).

\subsection{Test-Case Prioritisation}

The problem of choosing $\sigma_t$ is an instance of \emph{test-case
prioritisation}, studied for decades in regression testing.
\citet{rothermel2001prioritizing} formalised it as the search for a permutation
of the test pool that maximises how quickly faults are detected, and introduced
the two greedy families still in use:

\begin{itemize}[leftmargin=*,itemsep=3pt]
\item \textbf{Total} strategies rank a test case by a static property, usually
how many lines of the program it exercises.
\item \textbf{Additional} strategies rank by \emph{marginal} gain: how many
lines it exercises that nothing scheduled so far has exercised, re-ranking after
every pick. This family, written $\mathrm{RTP_{add}}$, is the stronger of the
two in most empirical comparisons and is used in this report as a classical
baseline.
\end{itemize}

\citet{elbaum2001cost} observed that the standard metric for prioritisation
assumes uniform test costs, and proposed a cost-cognizant replacement. That
paper is the ancestor of the cost term in this work. What it did not have is an
\emph{online} cost estimate: its costs are supplied before the run, whereas a
repair loop can measure them as it goes, for free, because it is running the
tests anyway.

\subsection{What \texorpdfstring{\frtp}{FRTP} Records and What It Discards}
\label{sec:bg-frtp}

\frtp~\citep{qi2013frtp} is the direct baseline for the scheduling half of this
work. It ranks each test case by its \emph{failing count}: the number of
candidate patches that test case has already rejected during this repair
session. The information is free, because the repair loop produces it as a
by-product, which is why \frtp's overhead is negligible.

Writing $k(x)$ for that count, \frtp\ schedules in decreasing $k$. Two things
are missing from $k$, and naming them precisely is what produces this work's
scheduling rule.

Figure~\ref{fig:frtp-gaps} makes both gaps concrete on four tests.

\begin{figure}[htbp]
\centering
\resizebox{0.95\textwidth}{!}{%
\begin{tikzpicture}[
  font=\small,
  lbl/.style={font=\small,anchor=west},
  hdr/.style={font=\small\bfseries,anchor=west},
  num/.style={font=\small,anchor=east},
  note/.style={font=\scriptsize\itshape,anchor=west,text=pendred,text width=44mm},
  grey/.style={font=\scriptsize\itshape,anchor=west,text=black!55},
  track/.style={draw=black!45,fill=black!6},
  kfill/.style={draw=bkblue!70,fill=bkblue!50},
  cost/.style={draw=black!45,fill=black!30}
]
\def\rw{0.09}   % units per run
\def\cs{0.90}   % units per CPU second

% ---- header ----
\node[hdr] at (0,0.85)  {test};
\node[hdr] at (1.3,0.85){runs ($k$ filled)};
\node[num,font=\small\bfseries] at (5.0,0.85){$k$};
\node[num,font=\small\bfseries] at (5.9,0.85){$n$};
\node[num,font=\small\bfseries] at (7.3,0.85){$K$};
\node[hdr] at (7.9,0.85){cost};
\draw[black!40] (0,0.58) -- (17.2,0.58);

% ---- T1 ----
\node[lbl] at (0,0.2) {$T_1$};
\draw[track] (1.3,0.05) rectangle ++({10*\rw},0.3);
\draw[kfill] (1.3,0.05) rectangle ++({8*\rw},0.3);
\node[num] at (5.0,0.2){8}; \node[num] at (5.9,0.2){10}; \node[num] at (7.3,0.2){0.80};
\draw[cost] (7.9,0.05) rectangle ++({0.05*\cs},0.3);
\node[lbl] at (9.4,0.2){0.05\,s};

% ---- T2 ----
\node[lbl] at (0,-0.7) {$T_2$};
\draw[track] (1.3,-0.85) rectangle ++({30*\rw},0.3);
\draw[kfill] (1.3,-0.85) rectangle ++({8*\rw},0.3);
\node[num] at (5.0,-0.7){8}; \node[num] at (5.9,-0.7){30}; \node[num] at (7.3,-0.7){0.27};
\draw[cost] (7.9,-0.85) rectangle ++({1.5*\cs},0.3);
\node[lbl] at (9.4,-0.7){1.50\,s};

% ---- T3 ----
\node[lbl] at (0,-1.6) {$T_3$};
\draw[track] (1.3,-1.75) rectangle ++({20*\rw},0.3);
\node[num] at (5.0,-1.6){0}; \node[num] at (5.9,-1.6){20}; \node[num] at (7.3,-1.6){0.00};
\draw[cost] (7.9,-1.75) rectangle ++({0.05*\cs},0.3);
\node[lbl] at (9.4,-1.6){0.05\,s};

% ---- T4 ----
\node[lbl] at (0,-2.5) {$T_4$};
\node[grey] at (1.3,-2.5){never run};
\node[num] at (5.0,-2.5){0}; \node[num] at (5.9,-2.5){0}; \node[num] at (7.3,-2.5){---};
\node[grey] at (7.9,-2.5){unknown};

% ---- the two gaps, bracketed on the right ----
\draw[pendred,thick] (11.2,0.2) -- (11.5,0.2) -- (11.5,-0.7) -- (11.2,-0.7);
\node[note] at (11.65,-0.25)
  {same $k$, so a count cannot order them --- yet $T_2$ rejects far less often
   \emph{and} costs $30\times$ more};
\draw[pendred,thick] (11.2,-1.6) -- (11.5,-1.6) -- (11.5,-2.5) -- (11.2,-2.5);
\node[note] at (11.65,-2.05)
  {same $k$ again --- but $T_3$ has $20$ runs of evidence and $T_4$ has none};

% ---- the three resulting orders ----
\draw[black!40] (0,-3.05) -- (17.2,-3.05);
\node[hdr] at (0,-3.5) {ranked by count $k$ \;(\frtp)};
\node[lbl] at (5.9,-3.5) {$\{T_1,T_2\}$ tied, then $\{T_3,T_4\}$ tied};
\node[note,text width=50mm] at (11.65,-3.5) {two ties, both hiding something};
\node[hdr] at (0,-4.15){ranked by rate $K$};
\node[lbl] at (5.9,-4.15){$T_1 \succ T_2 \succ T_4 \succ T_3$};
\node[grey] at (11.65,-4.15){ties broken; $T_3$ now ranks last};
\node[hdr] at (0,-4.8) {ranked by rate $\div$ cost};
\node[lbl] at (5.9,-4.8) {$T_1\,(16.0)\;\gg\;T_2\,(0.18)$};
\node[grey] at (11.65,-4.8){$T_2$ demoted by $90\times$};
\end{tikzpicture}}
\caption[Why a count of rejections is not enough]{The two things a rejection
count leaves out, on four tests. $T_1$ and $T_2$ have rejected the same number
of patches, so \frtp\ cannot order them --- yet $T_1$ rejects four candidates
in five and $T_2$ roughly one in four, and $T_2$ costs thirty times as much to
run. $T_3$ and $T_4$ also tie at zero, although $T_3$ has twenty runs of
evidence that it never rejects and $T_4$ has none. Reading the denominator
breaks both ties; dividing by cost then demotes the expensive test.}
\label{fig:frtp-gaps}
\end{figure}

\paragraph{The denominator is missing.} Let $n(x)$ be the number of times test
case $x$ has been \emph{run}. A count $k(x)$ without $n(x)$ cannot distinguish a
test case that was run ten times and rejected nothing from a test case that was
never run at all: both have $k = 0$. The ratio $k(x)/n(x)$ is the first quantity
that reads the evidence of \emph{not} rejecting, and it is also the first that
lets a test case whose early record was luck decay back towards the average.

\paragraph{The cost is missing.} Nothing in $k$ knows that one test runs in
milliseconds and another takes seconds. Under a uniform-cost assumption that is
harmless. When costs are spread out it is not, and it is worse than neutral if
expensive tests also tend to be more discriminating --- in that case ranking on
rejections alone actively steers the schedule towards the expensive end of the
pool.

\section{Repair Hypotheses and History in the Prompt}
\label{sec:bg-hypothesis}

\subsection{Patches Versus Hypotheses}

A \emph{repair hypothesis} is the idea behind a patch: what the model believes
is wrong and how it intends to fix it --- ``the loop bound is off by one'',
``integer division should be float division''. One hypothesis can be written
as many different patches, differing in variable names, statement order or
whitespace.

The distinction matters for measurement. A \emph{duplicate} patch, in the sense
used by \citet{chen2025conversational}, is a candidate whose normalised text
has already been generated this session. Two candidates that are textually
different but express the same hypothesis are not duplicates by that
definition, yet they are just as wasteful: both must be validated and both fail
for the same reason. The duplicate rate is therefore a \emph{lower bound} on
the redundancy a repair loop produces, which is the sense in which the reported
$65\,\%$ for \textsc{ChatRepair} is conservative.

\subsection{Why Showing the History Is Not Enough}

The natural fix for redundancy is to show the model what it has already tried.
Three published observations say this is necessary but not sufficient. A
transcript does not prevent repetition: \textsc{ChatRepair} keeps earlier
incorrect patches and their failures in the conversation, and
\citet{chen2025conversational} still measured $65\,\%$ of its patches to be
duplicates. Refinement preserves a bad hypothesis: when the underlying idea is
wrong, feedback-driven refinement produces superficial variants of it rather
than a new idea~\citep{le2026semantic}, and the loop pays for a validation that
was predictable in advance. And more context can hurt:
\citet{parasaram2025factselection} found repair effectiveness to be
\emph{non-monotonic} in the number of facts in the prompt, so past a point,
adding true and relevant information makes repair worse --- and a transcript
grows monotonically, so a long session necessarily walks past that point.

Taken together, the problem is not the \emph{absence} of history but its
\emph{form}. A growing transcript is the worst of both: it costs tokens every
round, it dilutes the few facts that matter, and it still leaves the model free
to re-propose.

\subsection{What a Hint Has to Be}

Two properties separate a useful hint from a transcript line.

\paragraph{It must be an observation, not a category.} A failure class ---
``the off-by-one family'' --- can only be phrased as a prohibition: \emph{do
not do that again}. A concrete test case can be phrased as a positive
requirement: \emph{on this input, produce exactly this output}. A requirement
tells the model what success looks like; a prohibition removes one option and
leaves it to guess the rest. This is also why the scheduler and the hint
contract are one system here: the scheduler's unit of memory is the test case,
and the test case is precisely the object that can be stated as a requirement.

\paragraph{It must be selected, not accumulated.} Because effectiveness is
non-monotonic in prompt content, the right number of hints is small and
roughly constant. Which hints those are changes from round to round, because
the history changes --- so selection must happen at every iteration, which is
what separates this from selecting static facts once before generation begins.

\section{Benchmark Background}
\label{sec:bg-benchmark}

ConDefects~\citep{wu2024condefects} collects real faulty submissions from
AtCoder programming contests, together with the author's later accepted version
and the contest's own test data: $1{,}254$ Java and $1{,}625$ Python faulty
programs. It was built to address data leakage in LLM evaluation, since its
faults post-date the training cut-offs of the models commonly benchmarked on
Defects4J~\citep{just2014defects4j}.

Four of its properties matter here. Every test input ships with its expected
output, so a rejection is a concrete input/output pair rather than a pass/fail
bit --- which is what lets a hint be stated as a positive requirement. A test
is one whole-program execution, so there is no call graph and nothing to
\emph{select away}; prioritisation is the only scheduling lever, which is the
structural reason the ``selection beats prioritisation'' finding of
\citet{lou2024aprregression} does not transfer here. Contest setters ship small
public samples alongside large adversarial inputs, so a single task's pool
spans a wide range of execution costs --- the spread a cost-aware schedule
needs. And a submission carries no problem statement, so every prompt must
supply worked examples, identically in every configuration, or the comparison
measures prompt content rather than policy.
